\documentclass[11pt]{article}

\usepackage[margin=25mm]{geometry}

\usepackage{amsmath,amssymb,fontspec,unicode-math}

\IfFontExistsTF{DejaVu Serif}{\setmainfont{DejaVu Serif}}{\setmainfont{Latin Modern Roman}}

\IfFontExistsTF{Latin Modern Math}{\setmathfont{Latin Modern Math}}{\setmathfont{DejaVu Math TeX Gyre}}

\usepackage{xurl,hyperref}

\setlength{\parindent}{0pt}\setlength{\parskip}{6pt}

\begin{document}

\title{RefLoRA: Refactored Low-Rank Adaptation for Efficient Fine-Tuning of Large Models}

\author{ELVIA preliminary AI analysis}\date{21 October 2026}\maketitle

Preliminary assessment: Major.\par

Independent scientific review remains pending. Author adoption has not been requested.

Theorem 4 reverses a negative-bound extremum and does not establish universally stronger balanced descent; the exact quadratic example confirms this failure. GD-to-AdamW scope and nonunique-surrogate branches remain local qualifications, with the balanced-factor formula preserved.

\section*{Review scope}

Latest-version, claim-focused review of the cited author-paper passages and applicable current author-release paths. The stated counterexamples and scope controls were checked analytically; no author program, old verifier, training run or benchmark was executed.

\section*{Problem 1: The balanced factor maximizes the negative lower bound}

Source: Theorem 4, p.6; Appendix B.5, p.20.

Assessment: supported; impact: major

\subsection*{Source claim (paraphrase)}

Theorem 4 says the balanced S minimizes a negative descent lower bound and uses it to explain improved descent.

\subsection*{Analysis}

The positive trace sum \(g(S)=\operatorname{tr}(XS)+\operatorname{tr}(YS^{-1})\) is minimized at the balanced \(S_b\). Multiplication by \(-\eta\|G\|_2^2\) reverses the extremum: the negative bound is maximized there, not minimized. In the admissible rank-one quadratic with \(A=B=e_1\), \(\eta=0.01\), the bound is \(-0.02\) at \(s=1\) and \(-0.0425\) at \(s=4\). The exact post-step losses from \(0.5\) are \(0.480298005\) and \(0.458498880\), so balancing need not give the larger one-step decrease.

\subsection*{Finding and scope}

Theorem4 reverses the extremum of its signed descent bound. The balanced factorization does not universally give the larger exact one-step loss decrease.

\subsection*{Proposed correction}

Correct the minimum to a maximum and remove or separately prove the stronger descent conclusion. The balanced-factor formula itself remains valid.

Proposed correction; independent verification and author adoption remain unestablished.

\section*{Problem 2: The GD invariance theorem does not cover AdamW}

Source: Theorem 6, p.7; Appendix B.6; reflora.py optimizer hooks.

Assessment: supported; impact: minor

\subsection*{Source claim (paraphrase)}

Theorem 6 proves a balanced gradient-descent factorization invariance.

\subsection*{Analysis}

Equivalent factors \((1,1)\) and \((2,1/2)\) give equal SGD weight change \(-0.0199\) in the scalar example, but initial AdamW normalization gives about \(-0.0199\) and \(-0.0249\) at learning rate \(0.01\). Coordinatewise adaptive normalization and optimizer state do not automatically respect the GD factor transformation. This is a scope boundary, not a counterexample to the GD theorem.

\subsection*{Finding and scope}

The invariance theorem concerns the specified gradient-descent update; current adaptive-optimizer integration does not inherit that property automatically. No contradiction to the restricted GD theorem is asserted.

\subsection*{Proposed correction}

State AdamW as a separate empirical algorithm. A general adaptive-state invariance proof remains unavailable.

Proposed correction; independent verification and author adoption remain unestablished.

\section*{Problem 3: Different surrogate optima need a consistent branch}

Source: Theorem 3, pp.4–5; linkage to Theorem 6 and default balanced rule.

Assessment: supported; impact: minor

\subsection*{Source claim (paraphrase)}

Theorem 3 admits multiple small-step surrogate minimizers, while Theorem 6 uses balanced factors.

\subsection*{Analysis}

For \(A=e_1,B=e_2,G=e_1e_1^\top,L=1,\eta=0.1\), the surrogate is \((s+1/s-10)^2\). Both \(s=5+\sqrt{24}\) and \(s=5-\sqrt{24}\) are global zero minima, but yield weight changes about \(-0.98989795e_1e_1^\top\) and \(-0.01010205e_1e_1^\top\). Arbitrary minimizer selection does not define a unique update. The default fixed balanced branch is not contradicted.

\subsection*{Finding and scope}

The optimizer-set statement does not select a unique update when different allowed surrogate minimizers are chosen. The default fixed balanced branch remains a countercontrol.

\subsection*{Proposed correction}

Specify the balanced or equivalence-consistent branch when combining the theorems; no default-algorithm error is inferred from nonuniqueness.

Proposed correction; independent verification and author adoption remain unestablished.

\section*{Limits}

Independent scientific adjudication remains pending.\par

No author code, model training or full benchmark evaluation was executed in this revision.\par

Version-specific source defects do not establish which code produced the published experiments.\par

Corrected benchmark scores, empirical effect sizes and author adoption remain unverified unless explicitly stated.\par

No new author-code execution, model spending, training, author contact or complete benchmark replication.\par

Historical numerical results are retained evidence rather than current reruns.\par

Local counterexamples do not invalidate every empirical result; missing historical run attribution remains unresolved.\par

\section*{Source identity}

arXiv:2505.18877v4 (current author version; updated 2026-05-29T23:09:57Z)

\url{https://arxiv.org/pdf/2505.18877v4}

Reviewed PDF SHA-256: \nolinkurl{e3ff630879d081cd46d8b8b0b62ecdb2b33e1672e90ca3423c9d532bc5f797c7}

The selected author/publisher source is the same work as the original case, matched by exact title, author identity and official version channel. The exact selected PDF SHA-256 was recomputed in the latest-version pass for Issue \#27. The PDF was downloaded from its version-specific official URL at the recorded time.

\end{document}
